\subsection{Data and Encoding}\label{app:data-specification}

The files in \path{reproducibility/data/} specify the coordinates of
$q$, the degree equations, the flag lists and the certificates for
\textbf{Lemmas~\ref{lem:finite-encoding} and~\ref{lem:rational-witness}}.

\subsubsection{Pattern Indices}\label{app:pattern-indices}

For each $n\in\{3,4,5,6\}$, the coordinates are indexed by the
relabeling orbits of the connected admissible patterns on $n$ positions.
As in \textbf{Lemma~\ref{lem:coordinates}}, the relabelings may take
either ordering of any marked pair to $(0,1)$.

List the unordered pairs of positions, written with the smaller
position first, in lexicographic order as
$e_0,e_1,\ldots,e_{\binom{n}{2}-1}$, with $e_0=01$. If $c_F(e_a)$ is
the color in $\{0,1,2\}$ assigned to the pair of positions $e_a$ by $F$, define
\begin{equation}\label{eq:pattern-code}
 \operatorname{code}(F)=\sum_{a=1}^{\binom{n}{2}-1}c_F(e_a)3^{a-1}.
\end{equation}
The root color $c_F(e_0)=2$ is fixed. For $n=3,4,5$, order the orbits by the lexicographically least tuple
$(c_F(e_1),\ldots,c_F(e_{\binom n2-1}))$ among the patterns $F$ in
each orbit. At order six, they are indexed in order of their least numerical
codes. \textbf{Table~\ref{tab:coordinates}} gives the file index ranges.
File index $j-1$ corresponds to $q_j$. The two-position pattern
$E_\star$ contributes the constant $1$.

\begin{table}[htbp]
\centering
\caption{Pattern counts and coordinate index ranges.}
\label{tab:coordinates}
\begin{tabular}{@{}rrrrr@{}}
\toprule
Order & Labeled patterns & Connected and admissible & Orbits & File indices\\
\midrule
3 & 9 & 8 & 5 & $0$--$4$\\
4 & 243 & 224 & 47 & $5$--$51$\\
5 & 19,683 & 18,576 & 676 & $52$--$727$\\
6 & 4,782,969 & 4,450,424 & 23,061 & $728$--$23,788$\\
\bottomrule
\end{tabular}
\end{table}

In the index file for the order of $F$, the entry at position
$\operatorname{code}(F)$ is $-1$ when $F$ is disconnected or
inadmissible. Otherwise, the files \path{idx3.bin}, \path{idx4.bin}
and \path{idx5.bin} store the file index from
\textbf{Table~\ref{tab:coordinates}}. For order six, the file index
is $728$ plus the nonnegative entry in
\path{order6_orbit_index.bin}. Each index file stores signed
32-bit little-endian integers. The files also number degree rows,
blocks and entries within blocks starting at zero.

\subsubsection{Degree Rows and Flag Lists}\label{app:flag-lists}

The file \path{degree.txt} gives the number of rows, followed by
the number $k$ of positions of the base pattern $F$, its code
$\operatorname{code}(F)$, and the position $u$ to which a new vertex is
attached.
For base orders $2,3,4,5$, the numbers of rows are $1,10,125,3380$,
respectively. For a row $(F,u)$ with base order $k$, the position $u$
specifies the pair $uk$ of color $1$ or $2$ in the extensions of
$F$ in \eqref{eq:degree-extension}. The coefficients count the
extensions by orbit, with $1$ subtracted from the base coordinate
when $F$ has more than two positions. The right-hand side is $1$
if $F=E_\star$ and $0$ otherwise.

The first line of \path{basis.txt} gives the number of coordinates and
the number of blocks. For each block, the file gives the order and code
of its type and the number of flags, followed by the order and code of
each flag in basis order. The first flag is the type itself. For each
listed flag order, the file contains one representative from each orbit
under permutations fixing the type positions.
\textbf{Table~\ref{tab:blocks}} gives the flag orders and sizes of the
bases.

\begin{table}[htbp]
\centering
\caption{The $57$ moment blocks and their $4{,}129$ flags.}
\label{tab:blocks}
\begin{tabular}{@{}lccr@{}}
\toprule
Type family & Flag orders & Block sizes & Blocks\\
\midrule
Marked two-position type & $2,3$ & $9$ & 1\\
All labeled three-position types & $3,4$ & $27$ & 8\\
Marked two-position type & $2,3,4$ & $133$ & 1\\
One type per four-position orbit & $4,5$ & $75,78,81$ & 47\\
\bottomrule
\end{tabular}
\end{table}

By \eqref{eq:moment-expansion}, each entry of $D_{\ell j}$ counts the
labeled completions of the corresponding pair of flags whose union
belongs to the orbit indexed by $j$. The two-position union
$E_\star$ contributes $1$ to the first diagonal entry of each
block with a two-position type. These are the only nonzero
entries of the constant matrices $K_\ell$.

The 57 blocks have 326,281 matrix entries in total. The coefficient
files represent both entries $(i,i')$ and $(i',i)$ when $i\ne i'$,
so the sum of entrywise products with a symmetric matrix counts both
off-diagonal entries in the trace product. The matrices $D_{\ell j}$
have $1{,}895{,}405$ nonzero entries in total.

\subsubsection{Sparse Coefficients and Certificates}

In the binary files, row and coordinate indices start at zero. Write
$N=23{,}789$ for the number of coordinates. For a degree row $a$ and
a coordinate index $j$, the key in
\path{degree_keys.bin} is $aN+j$. The corresponding signed coefficient
is in \path{degree_values.bin}. For the moment matrices, list the
blocks in \path{basis.txt} order, with the entries of each block
in row order. In a block of size $s$, entry $(i,i')$ has offset
$is+i'$ within that block.
For an entry at position $p$ in the full list and a coordinate index $j$,
the key in \path{keys.bin} is $pN+j$ and the corresponding coefficient
is in \path{values.bin}. Each coefficient is stored at the same
file position as its key. Keys are nonnegative and stored as signed
64-bit little-endian integers. Coefficients are signed 32-bit
little-endian integers. An absent key has coefficient zero.

The checker reconstructs the objective vector $o$ from
\path{idx4.bin} by adding $\frac{1}{2}$ to the coordinate of each
labeled order-four pattern with pair $23$ marked, as in
\eqref{eq:objective-vector}. The JSON certificate stores the
degree multiplier numerators and the integer vectors and weights in
\eqref{eq:integral-Q}, along with their denominators.
Each degree numerator is keyed by the base order, the full color
tuple, and the attachment position. An unlisted row has multiplier
zero.
The earlier certificate \path{original_sharp_certificate.json}
has $4{,}115$ positive-weight factors and gives the same bound
for the same degree equations and moment matrices. The file
\path{reproducibility/DATA_FORMAT.md} specifies the layouts and the
JSON fields of both certificates.
